\chapter{枚举与模式匹配}
\label{ch:10}

第 9 章的结构体解决的是「这些字段永远一起出现」。但有一类需求它搞不定：

\begin{tip}{提示}
一个值要么是这种，要么是那种。
\end{tip}

「一个形状」要么是圆（带半径），要么是方形（带边长）。「一次操作的结果」要么成功（带数据），要么失败（带原因）。

这种「多选一」的数据就是枚举。配合 \texttt{match}，它是 Magel 里最漂亮的一块。

\section{定义枚举}

\par
\medskip
\begin{lstlisting}[style=magel]
enum Message
{
    Quit,
    Move(int32 x, int32 y),
    Text(String text)
};
\end{lstlisting}
\par
\medskip

读法是：

\begin{itemize}

\item \texttt{Message} 这个类型，有三种变体（variant）：\texttt{Quit}、\texttt{Move}、\texttt{Text}；

\item \texttt{Quit} 不带数据；

\item \texttt{Move} 带两个 \texttt{int32}；

\item \texttt{Text} 带一个 \texttt{String}。

\end{itemize}

变体可以带字段，这在类型理论里叫「代数数据类型」——但用起来没那么玄：枚举就是「有几个可能形状，每个形状各自带自己的数据」。

枚举声明在 \texttt{\}} 处结束，尾分号可选，与 \texttt{struct} 和 \texttt{interface} 一致。

造一个值用 \texttt{类型:\allowbreak{}:\allowbreak{}变体}（片段，\texttt{Message} 假定已定义）：

\par
\medskip
\begin{lstlisting}[style=magel]
Message first = Message::Quit;
Message second = Message::Move((int32)3, (int32)7);
Message third = Message::Text("hello");
\end{lstlisting}
\par
\medskip

（\texttt{(\allowbreak{}int32)\allowbreak{}3} 那个强转是必须的——字面量默认是 \texttt{int64}，第 3 章的老规矩。）

\section{用 \texttt{match} 处理}

拿到一个枚举值，怎么知道它是哪个变体、里面的数据是什么？用 \texttt{match}：

\par
\medskip
\begin{lstlisting}[style=magel]
enum Message
{
    Quit,
    Move(int32 x, int32 y),
    Text(String text)
};

void show(Message message)
{
    match (message)
    {
        Message::Quit =>
        {
            println("quit");
        }
        Message::Move(x, y) =>
        {
            println("move {} {}", x, y);
        }
        Message::Text(text) =>
        {
            println("text {}", text);
        }
    }
}

int64 main()
{
    show(Message::Quit);
    show(Message::Move((int32)3, (int32)7));
    show(Message::Text("hello"));
    return 0;
}
\end{lstlisting}
\par
\medskip

\par
\medskip
\begin{lstlisting}[style=out]
quit
move 3 7
text hello
\end{lstlisting}
\par
\medskip

写法是：

\par
\medskip
\begin{lstlisting}[style=out]
match (要看的东西) {
    模式 => { 该做什么 }
    模式 => { 该做什么 }
}
\end{lstlisting}
\par
\medskip

注意三个语法细节：

\begin{enumerate}

\item 被匹配的表达式要用括号包起来：\texttt{match (\allowbreak{}message)\allowbreak{} \{}。

\item 箭头是 \texttt{=\allowbreak{}>\allowbreak{}}（和 C++ 的 \texttt{switch} 完全不同），每个分支的主体是一对花括号。

\item \texttt{match} 后面不加分号——它是语句；多写一个 \texttt{\};} 也能编译，但没必要。这一点和 \texttt{loop} 一样：\texttt{loop} 虽然是表达式，写在语句位置同样不补分号，只有被变量或别的表达式接走时，才按外层语句的规矩收尾。

\end{enumerate}

分支里的变量是「自动绑定」的。 \texttt{Message:\allowbreak{}:\allowbreak{}Move(\allowbreak{}x,\allowbreak{} y)\allowbreak{} =\allowbreak{}>\allowbreak{} .\allowbreak{}.\allowbreak{}.\allowbreak{}} 里写的 \texttt{x}、\texttt{y} 不是提前声明的变量，而是当场从变体里取出来的（片段，一条 \texttt{match} 分支）：

\par
\medskip
\begin{lstlisting}[style=magel]
Message::Move(x, y) =>
{
    println("move {} {}", x, y);
}
\end{lstlisting}
\par
\medskip

这一句同时干了两件事：判断「是不是 \texttt{Move}」，以及「把里面那两个 \texttt{int32} 取出来叫 \texttt{x}、\texttt{y}」。

\section{穷尽性：编译器帮你查漏}

这是 \texttt{match} 比 \texttt{if /\allowbreak{} else} 强得多的地方。

\par
\medskip
\begin{lstlisting}[style=magel]
enum Color
{
    Red, Green
};
\end{lstlisting}
\par
\medskip

\par
\medskip
\begin{lstlisting}[style=magel]
enum Color
{
    Red,
    Green
};

int64 name(Color color)
{
    match (color)
    {
        Color::Red =>
        {
            return 1;
        }
    }
    return 0;
}

int64 main()
{
    return 0;
}
\end{lstlisting}
\par
\medskip

\par
\medskip
\begin{lstlisting}[style=out]
magel: enum match 必须覆盖所有变体或提供 `_` 兜底分支；缺少：Color::Green
\end{lstlisting}
\par
\medskip

漏了一个变体是编译错误，不是运行期的意外。

这条规则的价值在于：以后你给枚举加一个变体，所有没处理它的 \texttt{match} 会立刻全部报错。编译器会挨个指出「你还没处理新来的那个」。

这才是「用类型保护重构」——改动带来的影响，编译器替你数清楚了。用 \texttt{if /\allowbreak{} else} 加字符串比较的写法，永远得不到这层保护。

实在想偷懒（比如只想处理其中一种），加一个兜底分支（\texttt{.\allowbreak{}.\allowbreak{}.\allowbreak{}} 是省略，不是真语法）：

\par
\medskip
\begin{lstlisting}[style=magel]
match (color)
{
    Color::Red =>
    {
        ...
    }
    _ =>
    {
        ...
    }
    // 剩下的都到这儿
}
\end{lstlisting}
\par
\medskip

\texttt{\_} 是通配符，表示「其他所有情况」。但每次写 \texttt{\_} 都是在放弃后面的检查——加新变体时不会再有提醒。能列全就列全。

\section{枚举方法}

枚举和结构体一样可以有方法，方法里用 \texttt{self} 指当前的枚举值：

\par
\medskip
\begin{lstlisting}[style=magel]
enum Shape
{
    Circle(int64 radius),
    Square(int64 side);

    int64 area()
    {
        match (self)
        {
            Shape::Circle(radius) =>
            {
                return 3 * radius * radius;
            }
            Shape::Square(side) =>
            {
                return side * side;
            }
        }
        return 0;
    }
};

int64 main()
{
    println("{}", Shape::Circle(2).area());
    return 0;
}
\end{lstlisting}
\par
\medskip

\par
\medskip
\begin{lstlisting}[style=out]
12
\end{lstlisting}
\par
\medskip

\texttt{match (\allowbreak{}self)\allowbreak{}} 就是「看看我自己是哪个变体」。把「一个形状的面积」这件事放在形状自己身上，用起来是 \texttt{shape.\allowbreak{}area(\allowbreak{})\allowbreak{}} ——这正是前面说的「把行为和数据放在一起」。

\subsection*{为枚举派生接口}

枚举也能使用 \texttt{\#derive}。相等判断先区分变体，再比较当前活动变体的载荷；哈希同时计入变体和载荷，克隆保留变体并克隆载荷。

\par
\medskip
\begin{lstlisting}[style=magel]
#derive(Equal, Hash, Clone)
enum BookPacket
{
    Empty,
    Text(String text)
}

int64 main()
{
    BookPacket first = BookPacket::Text("value");
    BookPacket second = first.clone();
    assert(!(first == BookPacket::Empty));
    println("{} {}", first == second, first.hash(first) == second.hash(second));
    return 0;
}
\end{lstlisting}
\par
\medskip

\par
\medskip
\begin{lstlisting}[style=out]
1 1
\end{lstlisting}
\par
\medskip

每个变体的载荷都必须满足所请求的接口；不能因为某个变体暂时没有在程序里构造，就跳过检查。与结构体一样，禁止派生 Drop 或 Closeable。包含无对应接口的结构体载荷时，应先为那个结构体实现或派生接口。

\section{不止能匹配枚举：标量也支持}

\texttt{match} 也能匹配整数、\texttt{char}、浮点、\texttt{bool} 和 String。整数与 \texttt{char} 还能使用区间模式：

\par
\medskip
\begin{lstlisting}[style=magel]
int64 grade(int64 score)
{
    match (score)
    {
        100 =>
        {
            return 3;
        }
        std::lang::rangeEq(90, 99) =>
        {
            return 2;
        }
        std::lang::rangeEq(60, 89) =>
        {
            return 1;
        }
        _ =>
        {
            return 0;
        }
    }
    return 0;
}

int64 main()
{
    println("{} {} {} {}", grade(100), grade(95), grade(60), grade(59));
    return 0;
}
\end{lstlisting}
\par
\medskip

\par
\medskip
\begin{lstlisting}[style=out]
3 2 1 0
\end{lstlisting}
\par
\medskip

区间模式统一使用标准库调用，可以显式写出类型参数：

\begin{center}
\begin{tabular}{ll}
\toprule
模式 & 命中条件 \\
\midrule
\texttt{std:\allowbreak{}:\allowbreak{}lang:\allowbreak{}:\allowbreak{}range<\allowbreak{}T>\allowbreak{}(\allowbreak{}a,\allowbreak{} b)\allowbreak{}} & \texttt{a <\allowbreak{}=\allowbreak{} x <\allowbreak{} b} \\
\texttt{std:\allowbreak{}:\allowbreak{}lang:\allowbreak{}:\allowbreak{}range\allowbreak{}Eq<\allowbreak{}T>\allowbreak{}(\allowbreak{}a,\allowbreak{} b)\allowbreak{}} & \texttt{a <\allowbreak{}=\allowbreak{} x <\allowbreak{}=\allowbreak{} b} \\
\bottomrule
\end{tabular}
\end{center}

区间匹配只支持整数和 \texttt{char}，不接受浮点、String 或自定义类型的区间。实现 \texttt{Equal<\allowbreak{}T>\allowbreak{}} 的自定义类型可以使用普通值模式，编译器通过相等协议比较。需要对浮点或 String 判断上下界时，在分支守卫中写出比较条件。

\subsection*{区间重叠的警告}

整数区间与区间、单值与区间的静态重叠会产生 \texttt{MGL1601} 警告，程序仍可编译。编译器按开闭端点和完整的 int64/uint64 精度判断相交，不会把警告改成类型错误，也不会重排分支；运行时仍按原顺序和守卫选择。重复的无守卫单值模式则沿用错误规则。命令行、JSON 与编辑器使用同一警告级别。

\section{守卫：先绑定，再加条件}

有时候光看「是哪个变体」不够，还要看「里面的值满足什么条件」。在模式后面加 \texttt{if}：

\par
\medskip
\begin{lstlisting}[style=magel]
enum Choice
{
    Value(String text)
};

int64 kind(Choice choice)
{
    match (choice)
    {
        Choice::Value(text) if text == "good" =>
        {
            return 1;
        }
        Choice::Value(text) =>
        {
            return 0;
        }
    }
    return -1;
}

int64 main()
{
    println("{} {}", kind(Choice::Value("good")), kind(Choice::Value("bad")));
    return 0;
}
\end{lstlisting}
\par
\medskip

\par
\medskip
\begin{lstlisting}[style=out]
1 0
\end{lstlisting}
\par
\medskip

\texttt{Choice:\allowbreak{}:\allowbreak{}Value(\allowbreak{}text)\allowbreak{} if text =\allowbreak{}=\allowbreak{} "good"} 的意思是：先确认是 \texttt{Value} 并取出 \texttt{text}，再判断 \texttt{text} 是不是 \texttt{"good"}。两个条件都满足才进这个分支。

注意第二条 \texttt{Choice:\allowbreak{}:\allowbreak{}Value(\allowbreak{}text)\allowbreak{} =\allowbreak{}>\allowbreak{} .\allowbreak{}.\allowbreak{}.\allowbreak{}} 不能省——\texttt{match} 还是要覆盖所有可能性，守卫不算「覆盖」。

\begin{tip}{提示}
守卫的关键字是 \texttt{if}。你在一些资料里可能看到写成 \texttt{when} 的版本——那份写法在当前工具链上是不认的（会报「期望 条件(if)」）。记住 \texttt{if} 就好。
\end{tip}

\section{什么时候用 \texttt{match}，什么时候用 \texttt{if}}

\begin{center}
\begin{tabular}{ll}
\toprule
场景 & 用什么 \\
\midrule
判断「一个值属于哪种形状」 & \texttt{match} —— 变体全列出来，编译器查漏 \\
一个数落在哪个区间 & \texttt{match} —— 比一串 \texttt{if /\allowbreak{} elif} 清楚 \\
复杂条件（多个变量、算出来的值） & \texttt{if /\allowbreak{} elif /\allowbreak{} else} \\
只有两个分支的简单判断 & \texttt{if /\allowbreak{} else} —— 别为了用 \texttt{match} 而用 \\
\bottomrule
\end{tabular}
\end{center}

一条经验：当「分支条件是某个值的形状/范围」时用 \texttt{match}；当「分支条件是几个东西的组合判断」时用 \texttt{if}。

\subsection*{\texttt{match} 表达式}

可以直接把匹配结果用于声明、赋值或调用。所有分支必须产生同类型值；枚举仍要求穷尽，守卫不算覆盖。

\par
\medskip
\begin{lstlisting}[style=magel]
enum Packet
{
    Data(String text), Empty
};
int64 main()
{
    Packet packet = Packet::Data("original");
    String label = match (packet)
    {
        Packet::Data(text) =>
        {
            text = "copy";
            yield text;
        }
        Packet::Empty =>
        {
            "empty"
        }
    };
    String original = match (packet)
    {
        Packet::Data(text) =>
        {
            text
        }
        Packet::Empty =>
        {
            "empty"
        }
    };
    println("{} {}", label, original);
    return 0;
}
\end{lstlisting}
\par
\medskip

\par
\medskip
\begin{lstlisting}[style=out]
copy original
\end{lstlisting}
\par
\medskip

表达式匹配的载荷绑定是值副本，修改副本不修改被匹配对象。单个结果写 \texttt{\{ value \}}，多语句块写 \texttt{yield value;}。外层声明或赋值仍需要分号。语句位置也支持 \texttt{match (\allowbreak{}.\allowbreak{}.\allowbreak{}.\allowbreak{})\allowbreak{} \{ .\allowbreak{}.\allowbreak{}.\allowbreak{} \};}；语句也允许省略末尾分号。

\section{小结}

\begin{center}
\begin{tabular}{ll}
\toprule
你会了 & 关键点 \\
\midrule
定义枚举 & \texttt{enum X \{ A,\allowbreak{} B(\allowbreak{}int64 v)\allowbreak{} \}} —— 结尾分号可选 \\
造值 & \texttt{X:\allowbreak{}:\allowbreak{}A}、\texttt{X:\allowbreak{}:\allowbreak{}B(\allowbreak{}3)\allowbreak{}} \\
\texttt{match} & 可作语句或表达式；表达式分支给出同类型值，外层声明仍加分号 \\
自动绑定 & \texttt{X:\allowbreak{}:\allowbreak{}B(\allowbreak{}v)\allowbreak{} =\allowbreak{}>\allowbreak{} .\allowbreak{}.\allowbreak{}.\allowbreak{}} 里的 \texttt{v} 是当场取出来的 \\
穷尽性 & 漏变体是编译错误；\texttt{\_} 兜底但会放弃检查 \\
标量匹配 & 普通值与整数 / char 区间：\texttt{100}、\texttt{std:\allowbreak{}:\allowbreak{}lang:\allowbreak{}:\allowbreak{}range\allowbreak{}Eq(\allowbreak{}90,\allowbreak{} 99)\allowbreak{}} \\
守卫 & \texttt{模式 if 条件 =\allowbreak{}>\allowbreak{} \{ .\allowbreak{}.\allowbreak{}.\allowbreak{} \}} \\
方法 & 枚举也能有方法，里面用 \texttt{match (\allowbreak{}self)\allowbreak{}} \\
\bottomrule
\end{tabular}
\end{center}

\section{下一步}

现在你能定义自己的类型了：结构体（几个字段一起）和枚举（几种形状选一）。但你还在用 \texttt{List<\allowbreak{}int64>\allowbreak{}}、\texttt{Map<\allowbreak{}String,\allowbreak{} int64>\allowbreak{}} 这些内建容器，而且总在数下标。

下一章把容器和遍历讲透。

翻到\hyperref[ch:11]{第 11 章 容器}。

\section{边界与验证}

枚举变体表达有限分类。match 的覆盖与守卫需要一起考虑；兜底分支可能隐藏新增变体，重叠警告也不应被误写成所有场合都拒绝编译。

\section{习题}

前两题用于理解和基本实现，第三题检查边界，最后一题为综合挑战。先独立完成，再核对参考解答。

\begin{enumerate}

\item 为任务状态设计三个枚举变体。

\item 用 match 为每个状态生成显示文字。

\item 增加一个失败状态，检查哪些匹配代码需要修改。

\item 挑战：设计带守卫的两个分支，解释书写顺序的影响。

\end{enumerate}

\subsection*{参考解答与验收要点}

\begin{enumerate}

\item 可用待办、进行中和完成；如果完成状态带信息，明确载荷类型。

\item 覆盖全部变体；带载荷时绑定后使用，不把判别值当任意整数。

\item 优先修改明确处理全部状态的位置；讨论兜底是否掩盖了需要单独处理的失败。

\item 给出同时满足两个条件的输入，说明实际选中分支，并检查不满足任何守卫时的后备路径。

\end{enumerate}

\section{番外：大灰狼漏掉的分支}

改编自传统童话《三只小猪》，以下为编程主题的独立改写。

大灰狼在山坡上开了一门拆屋生意。它的手册只有两条：草屋吹倒，木屋推倒。第一次接单是草屋，第二次是木屋，客户都给了好评。

第三只小猪交来一份砖屋委托，大灰狼翻到手册末页，那里写着“其余照办”。于是它先吹再推，忙到晚上，屋子仍在，自己的牙却松了一颗。

小猪问它是否承诺处理所有房屋。大灰狼说，前两次成功已经足够说明本领。小猪便把三种房屋摆在图上，让它指出每一种对应哪一个分支。

大灰狼想添一句“遇见麻烦便算木屋”，小猪不同意。它只好重新写下砖屋的处理结果：当前不能拆，必须明确返回。那天它少赚了一单，却终于知道“有一个兜底”与“认真处理了每一种情况”并不是同一回事。

\subsection*{技术对照}

枚举变体表达有限分类。match 的覆盖与守卫需要一起考虑；兜底分支可能隐藏新增变体，重叠警告也不应被误写成所有场合都拒绝编译。
